%------------------------
% English CV for Che CHENG
% Adapted from the COS Northeastern University Faculty LaTeX CV template
% Original template author: Zoe Kearney
% License: LaTeX Project Public License 1.3c
%------------------------

\documentclass[a4paper,9pt]{extarticle}

\usepackage[utf8]{inputenc}
\usepackage[T1]{fontenc}
\usepackage{lmodern}
\usepackage[a4paper,margin=0.68in]{geometry}
\usepackage{titlesec}
\usepackage{enumitem}
\usepackage{hyperref}
\usepackage{xcolor}

\definecolor{linkblue}{HTML}{1F5FBF}
\hypersetup{colorlinks=true,urlcolor=linkblue,linkcolor=linkblue}
\urlstyle{same}

\pagestyle{empty}
\setlength{\parindent}{0pt}
\setlength{\parskip}{0pt}
\setlength{\emergencystretch}{2em}
\setlist[itemize]{leftmargin=1.15em,itemsep=1pt,topsep=2pt,parsep=0pt}
\titleformat{\section}{\large\bfseries\uppercase}{}{0pt}{}[\titlerule]
\titlespacing*{\section}{0pt}{0.72\baselineskip}{0.38\baselineskip}

\newcommand{\role}[4]{%
  \textbf{#1}\hfill #2\\
  \textit{#3}\\
  \if\relax\detokenize{#4}\relax\else #4\fi
  \vspace{0.35\baselineskip}
}

\newcommand{\project}[4]{%
  \textbf{#1}\hfill \href{#4}{GitHub}\\
  #2\\
  \textit{Keywords: #3}
  \vspace{0.42\baselineskip}
}

\begin{document}
\small

\begin{center}
  {\Large\textbf{Che CHENG}}\\[3pt]
  MPhil Student, Computer and Information Engineering\\
  The Chinese University of Hong Kong, Shenzhen\\[2pt]
  Shenzhen, China \,|\, \href{mailto:checheng@link.cuhk.edu.cn}{checheng@link.cuhk.edu.cn}
  \,|\, \href{https://checheng117.github.io/}{Homepage}
  \,|\, \href{https://github.com/checheng117}{GitHub}
\end{center}

\section*{Profile}
\textbf{PhD-seeking MPhil student in embodied AI, robotic perception, and multimodal learning.}
My work focuses on structured 3D grounding, temporal interaction understanding, and
geometry-aware multimodal perception. I am interested in building reliable
perception-to-decision interfaces that expose intermediate state, uncertainty, and
diagnostic signals for embodied systems operating under occlusion, changing scene
state, and limited supervision.

\section*{Education}
\role
{The Chinese University of Hong Kong, Shenzhen}
{2025.09--2027.06}
{MPhil in Computer and Information Engineering}
{Research at iHEAR Lab on 3D vision, surgical automation, embodied intelligence,
multimodal learning, and robotic perception/control. Advisor:
\href{https://fxzhong.github.io/}{Prof. Fangxun Zhong}.}

\role
{Dalian Maritime University}
{2021.09--2025.06}
{Bachelor of Management in Big Data Management and Application}
{Minor in Software Engineering leading to a Bachelor of Engineering degree.}

\section*{Research Agenda}
\begin{itemize}
  \item \textbf{Structured 3D Grounding:} grounding objects, regions, spatial relations,
  and task context into structured representations for embodied perception.
  \item \textbf{Temporal Interaction Understanding:} modeling how agent-scene
  relationships evolve in dynamic, partially observable, and visually ambiguous scenes.
  \item \textbf{Reliable Embodied Decision Interfaces:} designing intermediate
  representations and evaluation protocols that make perception-driven decisions easier
  to inspect, diagnose, and improve.
\end{itemize}

\section*{Selected Open-Source Projects}
\textbf{Spatial Prism: Relation-Aware 3D Grounding for Embodied Perception}
\hfill \href{https://github.com/checheng117/relation-aware-3d-grounding}{GitHub}\\
\begin{itemize}
  \item Built a reproducible diagnostic project for relation-aware 3D visual
  grounding with scene-disjoint Nr3D/ReferIt3D-style evaluation, hard-subset
  analysis, coverage diagnostics, and controlled relation-scoring ablations.
  \item Preserved claim boundaries in the public report: reproduced baseline
  Acc@1/Acc@5, analyzed same-class clutter and high-clutter subsets, and reported
  controlled ablations without making state-of-the-art claims.
  \item Public artifacts include a course report, reproducibility notes, evidence map,
  diagnostic summaries, coverage reports, and aggregate tables; the README reports a
  4,255-sample scene-disjoint baseline and protocol-specific hard-subset diagnostics.
  \item The public summary separates reproduced baselines, frozen-logit diagnostics,
  trained controlled ablations, pilots, and negative findings to keep evidence levels
  auditable.
\end{itemize}
\textit{Keywords: 3D grounding; relation reasoning; embodied perception; diagnostics}
\vspace{0.34\baselineskip}

\textbf{RGB-D Spatial Grounding and Risk-Aware Robot Perception}
\hfill \href{https://github.com/checheng117/rgbd-spatial-grounding}{GitHub}\\
\begin{itemize}
  \item Implemented a real-time RGB-D pipeline for instance segmentation, target
  locking, 3D localization, point-cloud processing, object tracking, and
  nearest-distance risk estimation.
  \item Organized configurable scripts for Orbbec RGB-D capture, YOLO/SAM2-based
  segmentation, depth back-projection, mask alignment, offline sequence processing,
  and visualization reports with JSON metadata.
  \item The repository separates perception, risk estimation, calibration, utilities,
  control, tests, and configuration files; recorded sequences include RGB frames,
  depth frames, masks, camera intrinsics, and per-frame target-state records.
\end{itemize}
\textit{Keywords: RGB-D; point clouds; robot perception; risk estimation}
\vspace{0.34\baselineskip}

\textbf{EmbodiedSceneAgent: Structured Scene Memory and Replanning for Embodied Tasks}
\hfill \href{https://github.com/checheng117/EmbodiedSceneAgent}{GitHub}\\
\begin{itemize}
  \item Built a closed-loop embodied cognition codebase around observation, scene
  memory, planning, execution, verification, and replanning under changing scene
  state.
  \item Implemented structured planner contracts, parser/repair/validation stages,
  semantic acceptance filters, repeated-no-effect guards, and reproducible smoke
  and minimal-evaluation scripts.
  \item Public results report a 73-case diagnostic comparison with improvements in
  tool-use accuracy, target-match rate, and strict task proxy; the repository
  explicitly separates proxy metrics from official environment success rates.
  \item Public documentation includes project overview, showcase cases, limitations,
  architecture diagrams, setup scripts, and minimal verification commands for
  reproducing the smoke-test path.
  \item The README reports accepted revised-plan comparisons on a small controlled
  set while explicitly treating those cases as diagnostic rather than benchmark-scale
  generalization evidence.
\end{itemize}
\textit{Keywords: embodied agents; scene memory; verification; replanning}
\vspace{0.34\baselineskip}

\textbf{OpenGUI-RL: Verifiable GUI Grounding for Computer-Use Agents}
\hfill \href{https://github.com/checheng117/OpenGUI-RL}{GitHub}\\
\begin{itemize}
  \item Formulated instruction-conditioned GUI grounding as a scoped one-step
  contextual-bandit decision from screenshot and instruction to target element,
  click point, bounding box, and action type.
  \item Implemented candidate-aware supervised grounding, deterministic reward
  scoring, reward-labeled candidate generation, reranking analysis, and held-out
  transfer evaluation across Mind2Web, ScreenSpot-v2, and VisualWebBench.
  \item The public repository frames results as scoped benchmark evidence for GUI
  grounding, not as a claim about full browser automation or long-horizon agents.
  \item Public artifacts include quantitative summaries, tables, figures, data-format
  examples, reproduction notes, and explicit limitations on dataset redistribution,
  candidate protocols, and single-step evaluation scope.
  \item The public result summary highlights OCR/DOM candidate semantics, conditional
  reranking headroom, point-native transfer, and protocol sensitivity across GUI
  grounding benchmarks.
\end{itemize}
\textit{Keywords: GUI grounding; verifiable reward; computer-use agents; reranking}
\vspace{0.34\baselineskip}

\textbf{Verifiable Code Generation Post-Training for Open LLMs}
\hfill \href{https://github.com/checheng117/Verifiable-Code-Generation-Post-Training-Pipeline-for-Open-LLMs}{GitHub}\\
\begin{itemize}
  \item Developed an auditable post-training pipeline for verifiable code generation,
  centered on MBPP supervised fine-tuning, clean-view data construction, execution-
  based evaluation, and HumanEval external checking.
  \item Preserved frozen configs, run-stamped artifacts, data-card notes, non-
  fine-tuning baselines, training-cost audits, peak GPU memory audits, and a bounded
  GRPO exploration that is clearly separated from the main SFT-centered conclusion.
  \item Public archived runs include MBPP and HumanEval evaluation artifacts, clean-view
  SFT result rows, baseline comparisons, data preprocessing scripts, sandboxed
  execution checks, and assignment-facing frozen configuration files.
  \item The repository documents data licensing, processed split sizes, parser/evaluator
  risks, catastrophic-forgetting sanity checks, training-cost audits, and peak-memory
  profiling for the promoted runs.
\end{itemize}
\textit{Keywords: LLM post-training; code generation; SFT; execution-based evaluation}
\vspace{0.34\baselineskip}

\section*{Ongoing Research}
\begin{itemize}
  \item \textbf{Temporal interaction understanding in visually complex scenes:}
  studying how visual and motion cues can support reliable interpretation of
  fine-grained interactions under occlusion, ambiguity, and changing scene state.
  \item \textbf{Geometry-aware multimodal grounding for embodied perception:}
  exploring how task context, visual evidence, spatial structure, and uncertainty can
  support interpretable intermediate representations for reliable decision-making.
\end{itemize}
\textit{Technical details are limited until public manuscripts or preprints are available.}

\section*{Experience}
\textbf{Alibaba.com / Alibaba International Station}\hfill 2026.05--2026.09\\
\textit{AI Application Development Engineer Intern}
\vspace{0.35\baselineskip}

\textbf{ZTO Express Co., Ltd.}\hfill 2023.06--2023.09\\
\textit{Project-Based Intern}
\begin{itemize}
  \item Worked on operational data analysis for cross-belt sorting equipment, focusing
  on delayed fault recognition, load fluctuation, resource dispatching, and peak-period
  scheduling support.
  \item Used NumPy, Pandas, K-Means, ARIMA, and Hadoop-based workflows for data
  cleaning, anomaly detection, fault-pattern recognition, load forecasting, and
  peak-flow analysis; proposed a dynamic allocation strategy for tidal sorting chutes.
  \item Cleaned 10k+ equipment-operation records with an approximately 89\% cleaning
  rate; the project report recorded a 23\% sorting-efficiency improvement and about
  89\% peak-load forecasting accuracy.
  \item The resulting solution was later deployed at ZTO's Chongqing sorting center to
  support peak-period equipment-load forecasting and scheduling optimization.
\end{itemize}

\section*{Teaching}
\role
{Teaching Assistant}
{AY2025--2026, Term 2}
{CSC1006: Artificial Intelligence for Science and Engineering}
{The Chinese University of Hong Kong, Shenzhen.}

\section*{Technical Skills}
\begin{itemize}
  \item \textbf{Multimodal and agent systems:} VLMs, GUI agents, OCR/DOM grounding,
  structured spatial reasoning, grounded decision interfaces.
  \item \textbf{Embodied perception:} RGB-D perception, 3D scene memory, point-cloud
  processing, long-horizon planning, verification, replanning, failure analysis.
  \item \textbf{Training and evaluation:} PyTorch, PyTorch3D, OpenCV, Open3D,
  Transformers, PEFT/LoRA, instruction fine-tuning, reward-based reranking,
  reproducible evaluation pipelines.
  \item \textbf{Programming and data tooling:} Python-based research engineering,
  NumPy, Pandas, structured configuration, command-line experiment scripts,
  repository documentation, and LaTeX-based technical writing.
  \item \textbf{Vision and geometry workflows:} segmentation output integration,
  depth back-projection, point-cloud filtering, object-centric state tracking,
  spatial-distance computation, and qualitative visualization for debugging.
  \item \textbf{Experiment management:} smoke tests, ablation settings, saved run
  artifacts, error summaries, report tables, and reproducibility notes for public
  research-engineering repositories.
\end{itemize}

\section*{Awards and Honors}
\begin{itemize}
  \item Second Prize, China Undergraduate Mathematical Contest in Modeling, 2023.
  \item Silver Award, ``Internet+'' Innovation and Entrepreneurship Competition
  (Liaoning Division), 2023.
  \item Outstanding Graduate of Dalian City, 2025.
  \item Outstanding Student Scholarship and Comprehensive Excellence Scholarship,
  Dalian Maritime University, 2023 and 2024.
  \item Outstanding Student Leader / Merit Student / Outstanding League Member,
  Dalian Maritime University.
\end{itemize}

\section*{Last Updated}
May 2026.

\end{document}
